\subsection{Paths as Structured Filesystem References}

A path encodes parent chains, stems, suffixes, roots, and separator rules. Strings carry path text but expose only generic \texttt{str} methods; \texttt{Path} exposes filesystem-aware decomposition.

\subsubsection{Modern File System Manipulation: Object-Oriented Paths via \texttt{pathlib}}

\texttt{pathlib.Path} is an ergonomic structural abstraction: \texttt{/} overload composes components, \texttt{.name}/\texttt{.parent}/\texttt{.suffix} expose structure, and methods delegate to OS services.

Performance trade-off: every \texttt{Path(...)} construction, \texttt{/} join, or method call materializes reference-counted Python objects on the heap. High-volume path arithmetic in tight loops allocates more than primitive \texttt{os.path.join} on raw strings---acceptable for clarity in most code; hot paths may prefer string-based \texttt{os} routines.

\begin{verbatim}
from pathlib import Path

str_path = "data/chuck_facts.txt"
path_obj = Path("data") / "chuck_facts.txt"
print(str_path, path_obj.name, path_obj.suffix)
print(type(str_path), type(path_obj))
\end{verbatim}

\subsubsection{Absolute Paths, Relative Paths, and Current Working Directory Resolution}

Relative paths resolve against \texttt{Path.cwd()}; \texttt{resolve()} yields an absolute path. The same relative string can denote different files if the process working directory changes.

\begin{verbatim}
from pathlib import Path
rel = Path("data") / "answer.txt"
print(Path.cwd(), rel, rel.resolve())
\end{verbatim}

\subsubsection{Platform Separators: Windows Backslashes, POSIX Slashes, and Portable Path Construction}

Manual backslash concatenation breaks on escape sequences (\texttt{\textbackslash n}). \texttt{Path} composition avoids separator bugs; \texttt{PureWindowsPath}/\texttt{PurePosixPath} model syntax without touching the live filesystem.
